\section*{第 \ref{ch:b3:deformation-mechanisms} 章——变形机制与剪切带}

\begin{solution}{exo:b3:deformation-mechanisms:1}
镁橄榄石：$1573/2163 = 0.73$；冰：$263/273 = 0.96$。冰
远比镁橄榄石更接近其熔点，故更易蠕变：冰川
正是在地表流动。
\end{solution}

\begin{solution}{exo:b3:deformation-mechanisms:2}
$\gamma = \tan\ang{70} = 2.75$，$\tan\ang{80} = 5.67$。在 \ang{90} 附近，
正切值随角度变化极快：在 $\psi =
\ang{84}$ 处误差一度，$\gamma$ 就改变约五分之一。
\end{solution}

\begin{solution}{exo:b3:deformation-mechanisms:3}
$\sigma = (10^{3.56}/10)^{1/1.26} = 363.1^{0.794} = \qty{108}{MPa}$.
\end{solution}

\begin{solution}{exo:b3:deformation-mechanisms:4}
$E/RT = 223000/(8.314\times773.15) = 34.69$，故 $\dot\varepsilon =
\num{1.1e-4}\times20^4\times e^{-34.69} = \num{1.1e-4}\times\num{1.6e5}\times
\num{8.6e-16} = \qty{1.5e-14}{s^{-1}}$.
\end{solution}

\begin{solution}{exo:b3:deformation-mechanisms:5}
$\eta = \sigma/(2\dot\varepsilon)$：在 \qty{400}{\celsius} 时，$\num{65.4e6}/\num{2e-14} =
\qty{3.3e21}{Pa.s}$；在 \qty{500}{\celsius} 时，$\num{18.0e6}/\num{2e-14} =
\qty{9.0e20}{Pa.s}$。这些温度下，中地壳
的黏度与冰盖下方地幔的大致相当。
\end{solution}

\begin{solution}{exo:b3:deformation-mechanisms:6}
$n = \log(\num{1.6e-4}/\num{e-5})/\log2 = \log16/\log2 = 4$：\omterm{def:b3:deformation-mechanisms:creep}{位错蠕变}，
其\omterm{def:b3:deformation-mechanisms:flow-law}{应力指数}远大于 1。
\end{solution}

\begin{solution}{exo:b3:deformation-mechanisms:7}
$\ln10 = (E/R)(1/1000 - 1/1100)$，故 $E = 8.314\times2.3026/\num{9.091e-5} =
\qty{211}{kJ.mol^{-1}}$.
\end{solution}

\begin{solution}{exo:b3:deformation-mechanisms:8}
纳巴罗–赫林蠕变（$d^{-2}$）快 100 倍；科布尔蠕变（$d^{-3}$）快 1000 倍。
沿晶界扩散的原子所走路径，其截面 $\delta
d$ 随晶粒变小而缩小，此外还有距离更短、
需迁移的体积更小。
\end{solution}

\begin{solution}{exo:b3:deformation-mechanisms:9}
$\sigma_b^{2.5} = \num{3.59e-6}\,d^{-3}/\num{6.54e-17}$：当 $d = 1000$，
$\sigma_b = 54.9^{0.4} = \qty{5.0}{MPa}$；当 $d = 100$ 时，$\sigma_b =
\num{54900}^{0.4} = \qty{79}{MPa}$。斜率为 $-3/2.5 = -1.2$：晶粒细化十倍，
边界强度提高 $10^{1.2} = 16$ 倍。
\end{solution}

\begin{solution}{exo:b3:deformation-mechanisms:10}
两条路径中驱动浓度差相同，且
扩散距离均为 $d$；单位应变所需迁移的体积为 $d^3$。晶格
提供的截面 $\propto d^2$，故速率正比于 $d^2/(d\,d^3) =
d^{-2}$；晶界提供的截面 $\propto\delta d$，故 $\delta d/(d\,d^3) = \delta
d^{-3}$。晶界与晶格速率之比为 $(B_B\delta D_B)/(B_LD_Ld)$，
当 $d < (B_B/B_L)\,\delta D_B/D_L$ 时该比值大于 1。
\end{solution}

\begin{solution}{exo:b3:deformation-mechanisms:11}
$\eta \propto \dot\varepsilon^{(1-n)/n} = \dot\varepsilon^{-2.5/3.5}$：应变速率提高十倍，
则黏度降为原来的 $10^{0.714} = 5.2$ 分之一。恰好变形更快的条带
会变得更弱，从而承担更多变形：应变
向其中集中。
\end{solution}

\begin{solution}{exo:b3:deformation-mechanisms:12}
变化前：$\dot\varepsilon_{\mathrm{dis}} = \num{6.54e-17}\times10^{3.5} =
\qty{2.07e-13}{s^{-1}}$，$\dot\varepsilon_{\mathrm{dif}} = \num{3.59e-6}
\times10/3000^3 = \qty{1.3e-15}{s^{-1}}$：合计 \qty{2.08e-13}{s^{-1}}，
为\omterm{def:b3:deformation-mechanisms:creep}{位错蠕变}。变化后：$\dot\varepsilon_{\mathrm{dif}} = \num{3.59e-6}\times10/30^3 =
\qty{1.33e-9}{s^{-1}}$，为\omterm{def:b3:deformation-mechanisms:creep}{扩散蠕变}，在相同应力下快 \num{6400} 倍：
再结晶使该带大幅软化。
\end{solution}

\begin{solution}{pb:b3:deformation-mechanisms:1}
\textbf{1.} 岩脉与\omterm{def:b3:deformation-mechanisms:shear-zone}{剪切带}的夹角由 \ang{90} 变为 \ang{6}：$\psi =
\ang{84}$，$\gamma = \tan\ang{84} = 9.51$。

\textbf{2.} $\gamma w = 9.51\times400 = \qty{3804}{m}$，约为 \qty{3.8}{km}。

\textbf{3.} $S_1 = \sqrt{1 + 9.51^2/4} + 9.51/2 = 4.859 + 4.755 = 9.614$，
$S_2 = 1/S_1 = 0.104$，$S_1/S_2 = 92.4$。

\textbf{4.} $\tan2\theta' = 2/9.51 = 0.210$，$\theta' = \ang{5.9}$：
接近长轴的面理随 $\gamma$ 增大
而愈加逼近\omterm{def:b3:deformation-mechanisms:shear-zone}{剪切带}。

\textbf{5.} 沿岩脉向北追踪，岩脉在\omterm{def:b3:deformation-mechanisms:shear-zone}{剪切带}内被带向东，
离开\omterm{def:b3:deformation-mechanisms:shear-zone}{剪切带}后位置更靠东：北盘相对南盘向东移动。
从南盘望去，远侧向右移动：该\omterm{def:b3:deformation-mechanisms:shear-zone}{剪切带}为
右行。

\textbf{6.} 在包含拉伸线理的平面内（对走滑\omterm{def:b3:deformation-mechanisms:shear-zone}{剪切带}，线理此处为水平），
并垂直于竖直面理，即水平截面。其 S 面理
相对\omterm{def:b3:deformation-mechanisms:shear-zone}{剪切带}逆时针倾斜，朝向
东北东，与应变椭圆长轴的方向相似。

\textbf{7.} $\sigma = (10^{3.56}/20)^{1/1.26} = 181.5^{0.794} = \qty{62}{MPa}$。

\textbf{8.} $242.1^{0.794} = \qty{78}{MPa}$，对应 \qty{15}{\micro m}，
$145.2^{0.794} = \qty{52}{MPa}$，对应 \qty{25}{\micro m}。

\textbf{9.} 是的：实验的应力范围为 34 至 \qty{268}{MPa}。

\textbf{10.} 该\omterm{def:b3:deformation-mechanisms:shear-zone}{剪切带}在 \qty{450}{\celsius} 下变形，远低于实验的 800 至
\qty{1100}{\celsius}；与温度无关的关系可以推广
至此温度，
不过这仍是一种外推。

\textbf{11.} $D = 3631/100^{1.26} = 3631/331 = \qty{11}{\micro m}$.

\textbf{12.} $E/RT = 223000/(8.314\times723.15) = 37.09$，且 $e^{-37.09} =
\num{7.79e-17}$.

\textbf{13.} $\dot\varepsilon = \num{1.1e-4}\times62^4\times\num{7.79e-17} =
\num{1.1e-4}\times\num{1.478e7}\times\num{7.79e-17} = \qty{1.27e-13}{s^{-1}}$.

\textbf{14.} $\dot\gamma = 2\dot\varepsilon = \qty{2.53e-13}{s^{-1}}$.

\textbf{15.} $\eta = \num{62e6}/(2\times\num{1.27e-13}) = \qty{2.4e20}{Pa.s}$,
接近阿尔多湖下方上地幔的 \qty{3e20}{Pa.s}，该带
其强度与地幔相当。

\textbf{16.} 在 \qty{400}{\celsius} 时，$e^{-E/RT}$ 的变化因子为 0.064
（慢了十六倍）；在 \qty{500}{\celsius} 时为 11.0。

\textbf{17.} 快 $(78/62)^4 = 2.5$ 倍。

\textbf{18.} $t = 9.51/\num{2.53e-13} = \qty{3.75e13}{s}$.

\textbf{19.} $\num{3.75e13}/\num{3.156e10} = \qty{1190}{ka}$.

\textbf{20.} $\qty{3804}{m}/\qty{1.19e6}{\text{年}} = \qty{3.2}{mm/\text{年}}$.

\textbf{21.} 比主要板块边界慢十倍或更多，该带
为次要分支，或是若干共享同一边界运动的分支之一。

\textbf{22.} $\dot\gamma = 0.02/(\num{3.156e7}\times400) = \qty{1.58e-12}{s^{-1}}$,
快 6.2 倍；应力升高 $6.2^{1/4} = 1.58$ 倍，达到 \qty{98}{MPa}，
而晶粒尺寸将为 $3631/98^{1.26} = \qty{11}{\micro m}$。

\textbf{23.} 时间正比于 $\sigma^{-4}$：仅晶粒尺寸范围就使
结果为估计值的 $(62/78)^4 = 0.40$ 至 $(62/52)^4 = 2.0$ 倍，而
指数 $1.26 \pm 0.13$ 的不确定度，以及\omterm{def:b3:deformation-mechanisms:flow-law}{流动定律}常数的不确定度，
会进一步扩大其范围：结果只是一个数量级，而非确定的日期。

\textbf{24.} 瓦尔登\omterm{def:b3:deformation-mechanisms:shear-zone}{剪切带}积累其应变
大约用了 \textbf{\qty[text-series-to-math=true]{1190}{ka}}。
\end{solution}
